% !PRINTED 476
\weberpage{476}
\textbf{§\,129. \gesperrt{Anwendung der Transformation zweiter Ordnung zur
Berechnung von Klasseninvarianten.}}

Die Transformation zweiter Ordnung läßt sich, wenn nötig
in mehrmaliger Wiederholung, auf alle solche Diskriminanten
$-4m$ anwenden, für die $y^2 + m x^2$ eine Potenz von 2 ist, also
z.\,B. auf $m = 7$, $15$, $23$, $31$.

Diese Rechnungen sind aber beschwerlich, und wir werden
einfachere Wege finden, um in diesen Fällen zum Ziele zu kommen.
Bessere Dienste leistet die Transformation zweiter Ordnung, um
aus einer bekannten Klasseninvariante eine neue zu finden, die zu
einer Diskriminante gehört, die das Vierfache der ersteren ist.

Dazu führen folgende Formeln:
Nach § 34 ist:
\[
f_1(\omega)^8 + f_2(\omega)^8 = f(\omega)^8,
\]
\[
f_1(\omega)^8 f_2(\omega)^8 = \frac{16}{f(\omega)^8},
\]
woraus
\[
f_1(\omega)^8 - f_2(\omega)^8 = \frac{\sqrt{f(\omega)^{24} - 64}}{f(\omega)^4}.
\]
Das Vorzeichen der Wurzel wechselt nur für $f(\omega)^{24} = 64$,
also für $\omega = i$, und ist, solange $-i\omega$ reell und größer als 1
ist, positiv zu nehmen, da die linke Seite für ein verschwindendes
$q$ positiv unendlich wird. Es ergibt sich daraus:
\[
2 f_2(\omega)^8 = \frac{f(\omega)^{12} - \sqrt{f(\omega)^{24} - 64}}{f(\omega)^4},
\]
oder
\[
f_2(\omega)^8 f(\omega)^4 \left[f(\omega)^{12} + \sqrt{f(\omega)^{24} - 64}\right] = 32,
\]
und auf dieselbe Weise:
\[
f_2(\omega)^8 f_1(\omega)^4 \left[f_1(\omega)^{12} + \sqrt{f_1(\omega)^{24} + 64}\right] = 32.
\]
Ferner ist nach § 34, (16):
\[
f_1(2\omega)\,f_2(\omega) = \sqrt{2},
\]
woraus
\weq{1}{\begin{aligned}
2 f_1(2\omega)^8 &= f(\omega)^4\left[f(\omega)^{12} + \sqrt{f(\omega)^{24} - 64}\right],\\
&= f_1(\omega)^4\left[f_1(\omega)^{12} + \sqrt{f_1(\omega)^{24} + 64}\right].
\end{aligned}}
Diese Formeln können dazu dienen, wenn $f(\omega)$ oder $f_1(\omega)$
bekannt ist, $f_1(2\omega)$ zu berechnen, also $f_1(\sqrt{-4m})$ aus $f(\sqrt{-m})$
%CONT>>
